\section{Proposed biological validation protocol}
\label{sec:wetlab}
This section describes an experiment that could falsify the discovery
hypothesis. It has \emph{not} been performed; no simulated MIC or docking
output is reported as an observation.

\subsection{Peptide selection and processing}
The primary comparison is the 57-residue precursor of YdcA/P0ACW4 versus
the 37-residue mature segment after cleavage of the annotated 1--20 signal
peptide. The 20-residue signal segment alone is a necessary control because
the GNN score changes markedly after it is removed. Include the 57-residue
annealed derivative InstiAMP-09a-1a as a model-optimization control, plus
one sequence drawn without selection from the 852-protein screen as a
negative process control. A known antimicrobial reference peptide and a
vehicle-only well are assay controls. Because the mature segment contains
cysteines, synthesis and handling should specify oxidation state, purity,
and whether a disulfide arrangement is defined or a mixture; otherwise a
single primary-sequence prediction could map to several chemical species.

\subsection{Activity endpoint}
A standard broth-microdilution assay would span a twofold serial dilution
range appropriate to the synthesized material and use at least three
independent biological replicates. The endpoint is the lowest tested
concentration with no visible growth under the assay's prespecified
incubation and medium conditions. The experiment should record insolubility,
precipitation, and turbidity interference rather than counting them as
growth inhibition. The user-facing prediction is deliberately qualitative:
there might be a growth-inhibition effect, not a specific MIC value. A
continuous potency claim would require activity-labeled training data,
which the binary classifier does not use.

\subsection{Safety and specificity}
A mammalian-cell viability or hemolysis assay at overlapping concentrations
is necessary to distinguish selective antimicrobial action from nonspecific
membrane toxicity. Material preparation should be verified by mass
spectrometry and chromatography. These outcomes cannot be inferred from an
AMP/non-AMP classifier score. A 0.908 score is a ranking statistic, not a
validated 90.8\% chance of activity and not a therapeutic claim.

\subsection{Statistical decision rule}
Pre-register the primary comparison (mature YdcA versus vehicle) and the
replicate count before running the assay. Report individual replicate MICs,
not just an average; for twofold dilutions the endpoint is interval-censored,
so a log$_2$-scale analysis with an explicit censoring convention is more
honest than a decimal mean. A failure to inhibit growth at the highest
non-toxic, soluble concentration would directly weaken the model-generated
hypothesis. The analysis should not select one strain or one sequence after
looking at the activity panel without marking it exploratory.
